\begin{corollary}
Let $Q$ be a type and let $r:Q\to Q\to\mathsf{Prop}$, with the displayed
typeclass hypothesis $[\mathsf{IsPreorder}\ Q\ r]$.  Then the full tagged
hierarchy $\mathsf{PowerQ}(Q)$ equipped with the four-clause relation
$\mathsf{PowerRel}(r)$ preserves better-quasi-order:
\[
\mathsf{Bqo}(r)\longrightarrow
\mathsf{Bqo}(\mathsf{PowerRel}(r)).
\]
Here $\mathsf{PowerQ}(Q)$ has arbitrary nonempty small branching, as in
\noderef{PowerQ}, and is not restricted to countable nodes.  The four
clauses of $\mathsf{PowerRel}(r)$ compare atoms by $r$, compare an atom with
a node by an existential child condition, compare a node with an atom by a
universal child condition, and compare two nodes by a universal-existential
child condition.
\end{corollary}
\begin{proof}
Assume $\mathsf{Bqo}(r)$.  We prove
$\mathsf{Bqo}(\mathsf{PowerRel}(r))$ using the exact quantifiers in
\noderef{Bqo}.  Thus fix an arbitrary multi-sequence

\[
h:\mathsf{IncSeq}\to\mathsf{PowerQ}(Q)
\]

and assume that $h$ is locally constant.  To prove that it is not bad,
assume for contradiction
\[
\mathsf{BadMultiSequence}(\mathsf{PowerRel}(r))\,h.
\]
The hypotheses of \noderef{PowerQReflection} are now precisely available:
the displayed $[\mathsf{IsPreorder}\ Q\ r]$ is the preorder hypothesis,
and $h$ has the assumed local constancy and badness.  That proposition
therefore supplies a multi-sequence $g:\mathsf{IncSeq}\to Q$ such that
\[
\mathsf{LocallyConstant}(g),\qquad
\mathsf{BadMultiSequence}(r)\,g,
\qquad
\forall x\in\mathsf{IncSeq},\quad
g(x)\in\mathsf{PowerSupport}(h(x)).
\]
In particular, the second component says that $g$ is bad for $r$ and the
first says that it is locally constant.  But the definition in
\noderef{Bqo}, applied to the assumed $\mathsf{Bqo}(r)$ and this particular
$g$, says
\[
\neg\mathsf{BadMultiSequence}(r)\,g,
\]
contradicting the second component above.  Hence the arbitrary locally
constant $h$ is not bad, and the defining quantifiers in \noderef{Bqo}
give $\mathsf{Bqo}(\mathsf{PowerRel}(r))$.
\end{proof}
